\chapter{Tán xạ lần cuối và sự phát sinh baryon}

\section{Nhìn ra xa là nhìn về quá khứ}

Ở những khoảng cách khác nhau, ta nhìn thấy vũ trụ ở những thời điểm khác nhau. Ở gần, thời gian nhìn lùi không quá quan trọng. Vũ trụ của một năm trước, hay một trăm năm trước, không khác vũ trụ hiện nay lắm. Nhưng nếu nhìn đủ xa, ta thấy vũ trụ trẻ hơn. Vũ trụ trẻ hơn thì đậm đặc hơn, và vũ trụ đậm đặc hơn thì nóng hơn.

Ở một khoảng cách đủ lớn nào đó, vũ trụ nóng xấp xỉ bề mặt Mặt Trời. Nó không có vẻ như vậy đối với ta vì các photon đã bị dịch đỏ trên đường tới đây. Hệ số dịch đỏ cỡ một nghìn có nghĩa là nhiệt độ quan sát được thấp hơn đúng cùng hệ số:
\[
T_{\mathrm{obs}} \sim \frac{T_{\mathrm{emit}}}{1+z},
\qquad
1+z \sim 10^3.
\]\footnote{Biểu thức hiển thị này là bản chép biên tập cô đọng của lập luận bằng lời về dịch đỏ tại 00:02:20--00:02:39.}

Đó là nơi vũ trụ trở nên mờ đục. Vũ trụ chủ yếu là hiđrô, cùng một ít heli. Trên nhiệt độ vài nghìn độ, hiđrô bị ion hóa. Vì vậy ở phía bên kia bề mặt đó không có các nguyên tử trung hòa mà là các ion: chủ yếu là proton và electron chuyển động tự do. Các photon ở đó cũng giàu năng lượng hơn, về cơ bản là tia X. Nếu một nguyên tử tình cờ hình thành, nó sẽ nhanh chóng bị ion hóa trở lại do va chạm hoặc do một photon.

Một khí ion hóa là một chất dẫn điện, và chất dẫn không cho bức xạ truyền qua một cách thông suốt. Các điện tích đáp ứng với điện trường, dao động qua lại, và phát xạ lại bức xạ. Vì thế vượt quá khoảng cách đó thì vũ trụ là mờ đục đối với ta. Bề mặt tán xạ lần cuối giống như bề mặt nhìn thấy của Mặt Trời: không phải một bức tường vật chất, mà là nơi môi trường chuyển khá đột ngột từ mờ đục sang trong suốt.

Ở đó, photon áp đảo về số lượng, với hệ số khoảng \(10^8\). Ngày nay cũng vậy, dù ngày nay các photon có năng lượng thấp đến mức chúng không phải là một phần quan trọng của mật độ năng lượng hiện tại. Một khi nhiệt độ hạ xuống dưới nhiệt độ ion hóa của hiđrô, vũ trụ trở nên trong suốt và ta có thể nhìn xuyên qua nó.

Tầm nhìn trực tiếp chấm dứt ở đó. Câu hỏi tiếp theo là vật lý nguyên tử, hạt nhân và hạt đã biết còn có thể đưa ta lùi xa thêm bao nhiêu về quá khứ.

\section{Lùi xa hơn nữa về quá khứ}

Ngay phía trong bề mặt tán xạ lần cuối, ngoài photon ra, các hạt chính là proton và electron, cùng một số hạt nhân heli và vài deuteron. Hầu như không có phản proton và hầu như không có positron. Nhưng lùi xa hơn nữa, nhiệt độ rốt cuộc trở nên đủ cao để hai photon có thể va chạm và tạo ra một electron và một positron:
\[
\gamma + \gamma \leftrightarrow e^- + e^+.
\]

Ngưỡng cơ bản là các photon ít nhất phải cung cấp năng lượng nghỉ của cặp đó:
\[
E \ge 2m_e c^2.
\]
Dưới ngưỡng ấy, ánh sáng thông thường không làm được việc này. Trên ngưỡng ấy, các cặp electron--positron có thể được tạo ra, và phản ứng ngược cũng có thể xảy ra. Một khối vật chất nóng ở nhiệt độ đó sẽ chứa các phản ứng hạt đóng vai trò mà các phản ứng hóa học đóng trong vật chất thông thường.

Ngay phía trên ngưỡng vẫn có nhiều photon hơn electron và positron. Khối lượng nghỉ của electron rất tốn năng lượng. Nhưng khi nhiệt độ tăng lên cao hơn ngưỡng rất nhiều, năng lượng nghỉ của electron trở nên kém quan trọng hơn năng lượng động học, và các mật độ trở nên tương đương. Khi đó electron, positron và photon đều có gần như cùng mật độ và va chạm loạn xạ dữ dội.

\begin{figure}[ht]
\centering
\includegraphics[width=0.62\textwidth]{lecture_06_figure_02.png}
\caption{Phác họa ngưỡng tạo cặp so sánh các thang khối lượng của electron và proton.}
\lectureframe{00:18:47}
\end{figure}

Câu hỏi tiếp theo là proton sẽ ra sao. Chúng đã hiện diện sẵn, và đến lúc ấy mật độ của chúng rất cao, nhưng vẫn chưa có đủ năng lượng nhiệt để tạo ra dồi dào các cặp proton--phản proton. Lý do rất đơn giản:
\[
m_p \approx 2000\,m_e,
\qquad
E \ge 2m_p c^2
\]
đó là một ngưỡng cao hơn nhiều.\editorialnote{Trong khung hình, nhãn khối lượng thấp ở phía dưới kém rõ hơn nhãn \(2m_e c^2\) ở phía trên, nhưng phép so sánh được phát biểu bằng lời một cách tường minh là giữa các ngưỡng của electron và proton.}

Mức này đã nóng hơn rất nhiều so với bên trong một ngôi sao, nơi phản hạt không hiện diện với độ phong phú đáng kể nào.

Nhưng ngay cả mô tả đó cũng chỉ là tạm thời. Đến khi nhiệt độ đủ cao để tạo ra proton và phản proton một cách tự do, thì nó cũng đủ cao để đập vỡ proton thành các quark. Vì vậy hỗn hợp cực nóng ấy được mô tả đúng hơn là một hỗn hợp gồm electron, positron, quark, phản quark, photon, neutrino, gluon và các loài hạt khác. Một phản ứng sơ đồ đơn giản trên bảng là
\[
e^- + e^+ \to \gamma \to q + \bar q.
\]

\begin{classroomqa}
\lecturetimestamp{00:20:05}
\audiencequestion{Là meson hay tiền-quark?}
\lecturerresponse{Meson được tạo thành từ quark. Một khi nhiệt độ đủ cao để phá vỡ một proton, nó cũng phá vỡ các meson.}
\end{classroomqa}

Cho tới đây, lập luận đã đi ngược về thời gian. Bây giờ hãy đảo chiều nó lại và hỏi điều gì xảy ra khi hỗn hợp nóng đối xứng nguội đi.

\section{Tính trung hòa điện và một vũ trụ nóng đối xứng}

Ở nhiệt độ đủ cao, vũ trụ chứa mật độ khổng lồ của hạt và phản hạt. Electron và positron xuất hiện theo cặp, quark và phản quark xuất hiện theo cặp, và bức tranh tự nhiên là độ phong phú gần như bằng nhau của mỗi hạt và phản hạt tương ứng của nó.

Tính trung hòa điện cần được hiểu cẩn thận đôi chút. Ngày nay vũ trụ là trung hòa điện. Một mẹo ghi nhớ hữu ích là hình dung, không phải vũ trụ thực của chúng ta, mà là một vũ trụ khép kín và hữu hạn, như một 3-cầu. Trước hết hãy nghĩ tới một mặt cầu hai chiều như một mô hình. Các đường sức điện không thể bắt đầu hay kết thúc trừ ở điện tích. Điện tích dương phát ra các đường sức; điện tích âm hấp thụ chúng. Trên một mặt cầu khép kín, các đường ấy không thể đơn giản thoát ra vô cực. Chúng hội tụ lại ở phía đối diện. Vì vậy, như một mẹo ghi nhớ, một vũ trụ khép kín và hữu hạn không thể mang điện tích thuần. Việc vũ trụ thực của chúng ta có thật sự khép kín hay không không phải là điều cốt yếu. Điều cốt yếu là ở quy mô lớn, các điện tích dương và âm dường như triệt tiêu nhau một cách cực kỳ tốt.

Tính trung hòa điện \emph{không} có nghĩa là số proton bằng số phản proton. Nó có nghĩa là tổng điện tích dương bằng tổng điện tích âm. Proton, positron và mọi loài hạt mang điện dương khác gộp lại cân bằng với electron và mọi loài hạt mang điện âm khác gộp lại.

Ngược về quá khứ xa đến mức vật lý đã biết còn có thể được tin cậy, và có lẽ còn xa hơn một chút nữa, vũ trụ đã đủ nóng để duy trì dồi dào cả hạt lẫn phản hạt. Bất kỳ thí nghiệm máy gia tốc nào ta có thể nêu tên cũng đều có một tương tự của nó trong vũ trụ sơ khai ở một giai đoạn nào đó. Giả thuyết tự nhiên là một khởi đầu gần đối xứng: số electron bằng số positron, số quark bằng số phản quark, và, nếu dùng proton và phản proton làm đại diện cho baryon và phản baryon,
\[
N_p = N_{\bar p}.
\]

Giả thuyết ấy hấp dẫn vì các định luật vật lý gần như đối xứng giữa hạt và phản hạt. Một positron ứng xử rất giống một electron, không hoàn toàn nhưng rất gần như thế. Chính sự gần đối xứng đó là vấn đề.

\section{Cho bộ phim chạy theo chiều thuận}

Giả sử vũ trụ bắt đầu trong một trạng thái nóng với độ phong phú bằng nhau hoặc gần bằng nhau của hạt và phản hạt, rồi nguội dần.

Trong khi nhiệt độ còn đủ cao, proton và phản proton có thể được tạo ra từ các va chạm năng lượng cao, và chúng có thể hủy nhau trở lại thành photon. Nhưng rồi cuối cùng nhiệt độ hạ xuống dưới ngưỡng proton. Khi đó các cặp proton--phản proton vẫn có thể hủy nhau thành photon, nhưng sau một chút giãn nở, các photon ấy đã mất đủ năng lượng nên không còn có thể biến ngược trở lại thành các cặp proton--phản proton nữa. Đây là con đường một chiều. Về sau điều tương tự xảy ra với electron và positron, ở một nhiệt độ thấp hơn ngưỡng proton rất nhiều nhưng vẫn còn cao hơn xa nhiệt độ trong sao.

Bởi vì sự giãn nở chậm so với các tốc độ phản ứng vi mô, vũ trụ phần lớn vẫn ở gần cân bằng nhiệt trong khi việc này diễn ra. Đó là điều ở đây được hiểu bởi sự giãn nở đoạn nhiệt. Trong bối cảnh ấy, phát biểu chính xác hơn là entropy thay đổi rất ít. Đối với một hỗn hợp hạt nóng, entropy về cơ bản chính là số hạt, sai khác chỉ bởi các hệ số cỡ đơn vị. Vì thế proton, phản proton, electron và positron bị hủy đi được thay thế xấp xỉ bằng photon.

Nếu số proton và phản proton chính xác bằng nhau, chúng sẽ đơn giản hủy nhau hết. Điều tương tự sẽ xảy ra với electron và positron. Khi đó chỉ còn lại photon. Nhưng nếu proton chỉ dư ra so với phản proton dù rất ít, thì các cặp ghép tương ứng sẽ hủy nhau và phần dư không ghép cặp sẽ còn lại. Một ví dụ cụ thể cho thấy điều đó: nếu có một trăm proton và năm mươi phản proton, thì năm mươi phản proton sẽ tìm thấy năm mươi proton, hủy nhau với chúng, và năm mươi proton sẽ còn lại. Khi đó, trung hòa điện đòi hỏi một phần dư electron so với positron cùng kiểu như vậy.

Điều này cho phép suy ra độ bất đối xứng thời kỳ đầu từ số photon hiện nay. Ngày nay số photon lớn hơn số proton hoặc electron khoảng \(10^9\) đến \(10^{10}\) lần:
\[
\frac{N_\gamma}{N_p} \sim 10^9\text{--}10^{10}.
\]
Vì vậy trong hỗn hợp nóng thuở đầu chỉ có khoảng 1 proton dư ra trên mỗi tỷ cặp proton--phản proton. Sự lệch cực nhỏ đó là toàn bộ lý do khiến nguyên tử vẫn còn tồn tại.

\begin{classroomqa}
\lecturetimestamp{00:42:25}
\audiencequestion{Tại sao không xem phần dư này là một hiệu ứng thống kê, được tạo ra trong quá trình nguội đi với một xác suất rất nhỏ?}
\lecturerresponse{Để có thêm một proton từ một trạng thái đối xứng ban đầu, ta sẽ phải tạo ra một proton mà không đồng thời tạo ra một phản proton. Đồng thời cũng phải tạo ra cùng với nó điện tích âm, chẳng hạn một electron. Mô hình chuẩn đã biết không chứa quá trình năng lượng thấp nào như vậy. Ở những năng lượng đã được kiểm nghiệm đó, số baryon được bảo toàn, nên sự nguội đi thông thường không sinh ra phần dư ròng ấy.}
\end{classroomqa}

Vì vậy điện tích không phải là trở ngại. Một ví dụ sơ đồ sẽ là
\[
\gamma + \gamma \to p + e^-.
\]\footnote{Biểu thức này là bản chép biên tập của ví dụ bằng lời tại 00:44:16--00:46:35. Mục đích của nó là tách riêng điểm mấu chốt rằng điện tích có thể được bảo toàn trong khi số baryon bị vi phạm.}
Xét tổng thể thì quá trình này trung hòa điện. Nhưng nó vi phạm một quy luật thực nghiệm khác. Trong mô hình chuẩn đã được kiểm nghiệm, người ta không thể tạo ra một proton mà không đồng thời tạo ra một phản proton. Với mục đích hiện tại, đại lượng bảo toàn liên quan là số baryon,
\[
B=(N_p+N_n)-(N_{\bar p}+N_{\bar n}).
\]\footnote{Công thức gọn này hình thức hóa định nghĩa bằng lời được nêu tại 00:45:18--00:45:40.}
Ở các năng lượng đã biết, đại lượng đó được bảo toàn, nên độ bất đối xứng không thể được tạo ra trong phần muộn, đã được kiểm nghiệm kỹ, của lịch sử nguội đi.

\section{Phát sinh baryon}

Đây là chủ đề gọi là sự phát sinh baryon. Giả sử vũ trụ thực sự đã bắt đầu với số baryon và phản baryon bằng nhau. Các định luật vật lý sẽ phải cho phép điều gì để làm lệch trạng thái đối xứng đó thành phần dư nhỏ bé đã được quan sát?

Trên bảng, điều này được rút gọn thành hai biểu thức ký hiệu,
\[
N_p=N_{\bar p},
\qquad
N_p-N_{\bar p},
\]
rồi tiếp đó thành một danh sách các điều kiện liên hệ biểu thức thứ nhất với biểu thức thứ hai.

\begin{figure}[ht]
\centering
\includegraphics[width=0.72\textwidth]{lecture_06_figure_04.png}
\caption{Các điều kiện để tạo ra một bất đối xứng baryon từ một trạng thái đối xứng ban đầu.}
\lectureframe{01:22:26}
\end{figure}

\editorialnote{Trong khung hình này, các từ trong danh sách đánh số chỉ đọc được một phần. Ba mục dưới đây bám theo phần giải thích bằng lời đi kèm.}

Điều kiện thứ nhất là sự vi phạm số baryon. Nếu số baryon là một định luật bảo toàn chính xác, thì một vũ trụ khởi đầu với số baryon ròng bằng không sẽ vẫn như vậy.

Điều kiện thứ hai là một sự thiên lệch giữa hạt và phản hạt. Chỉ riêng sự vi phạm số baryon là chưa đủ. Nếu tự nhiên cho phép một quá trình tạo proton và quá trình phản chiếu tương ứng tạo phản proton với xác suất như nhau, thì độ mất cân bằng trung bình vẫn bằng không. Cần các tốc độ xảy ra khác nhau.

Các dao động thống kê thông thường cũng quá nhỏ. Nếu tổng số hạt, kể cả photon, vào khoảng \(10^{90}\), thì một dao động cỡ \(\sqrt{N}\) chỉ cho khoảng \(10^{45}\), hoàn toàn không đủ để giải thích xấp xỉ \(10^{80}\) proton trong vũ trụ.

Điều kiện thứ ba là vũ trụ phải, vào một thời điểm nào đó, đủ xa khỏi cân bằng. Ở đây lập luận dùng định lý CPT. Nếu một khí nóng đã ở cân bằng mãi mãi, thì một đoạn phim chạy ngược trông giống hệt một đoạn phim chạy xuôi. Trong bối cảnh đó, đối xứng giữa hạt và phản hạt được khôi phục một cách hiệu dụng bởi đối xứng chính xác kết hợp liên hợp điện tích, phép chẵn lẻ và đảo ngược thời gian. Muốn có một thiên lệch thực sự được đóng băng lại, lịch sử thực tế phải trông khác nhau khi đi xuôi và đi ngược trên các thang thời gian vi mô. Ở đủ sớm, sự giãn nở của vũ trụ đủ nhanh để làm được điều đó.

Vì vậy ba điều kiện là:
\begin{enumerate}
\item sự vi phạm số baryon;
\item sự vi phạm đối xứng giữa hạt và phản hạt, chính xác hơn là vi phạm CP;
\item một giai đoạn ngoài cân bằng nhiệt.
\end{enumerate}

\begin{classroomqa}
\lecturetimestamp{01:09:08}
\audiencequestion{Nếu ghép các điều kiện này với CPT, thì điều kiện 2 thực sự là CP còn điều kiện 3 thực sự là vấn đề chiều thời gian, phải không?}
\lecturerresponse{Đúng. Điều kiện 2 thực ra nên được phát biểu là vi phạm CP. Điều kiện 3 là việc lịch sử thực tế ở ngoài cân bằng, không phải một định luật bảo toàn riêng biệt và cũng không phải cùng một chuyện với vi phạm số baryon.}
\end{classroomqa}

Với ba điều kiện này, ngay cả một vũ trụ khởi đầu với số proton và phản proton bằng nhau cũng có thể bị làm lệch theo hướng này hay hướng kia. Điều còn chưa chắc chắn là liệu tự nhiên có thực sự vi phạm số baryon hay không, và nếu có thì bằng cách nào.

\section{Điều đã biết, và điều chưa biết}

Phần còn bỏ ngỏ của câu chuyện là sự vi phạm số baryon. Chúng ta chưa bao giờ trực tiếp quan sát thấy trong tự nhiên một phản ứng làm thay đổi
\[
(N_p+N_n)-(N_{\bar p}+N_{\bar n}).
\]
Điều đó có thể chỉ có nghĩa là các phản ứng liên quan đòi hỏi nhiệt độ cao hơn rất nhiều so với bất cứ điều gì đã được kiểm nghiệm trực tiếp trong phòng thí nghiệm, hoặc rằng chúng cực kỳ khó xảy ra ở năng lượng thấp. Nhưng hiện nay đây vẫn là dấu hỏi.

Đó là lý do vì sao phân rã proton lại quan trọng đến vậy trong cuộc thảo luận. Nếu số baryon không phải là chính xác, thì về nguyên tắc proton phải không bền. Nhưng nó lại bền một cách đáng kinh ngạc. Cận dưới được nêu ở đây có cỡ
\[
\tau_p \gtrsim 10^{34}\ \text{years}.
\]
Điều này đến từ những thí nghiệm khổng lồ dùng các thể tích nước rất lớn, được gắn đầu dò ánh sáng và đặt sâu dưới lòng đất. Nó không bác bỏ sự vi phạm số baryon. Nó chỉ có nghĩa là bất kỳ sự vi phạm nào như vậy đều phải bị kìm hãm ở mức phi thường tại năng lượng thấp.

Một nhận xét khác thuộc vật lý đã biết cũng đi vào cuộc thảo luận. Quanh \(100\) đến \(250\,\mathrm{GeV}\), các tương tác yếu trở nên giống các tương tác điện từ hơn nhiều. Điều đó có thể quan trọng trong các mô hình chi tiết, nhưng không phải là bước đột ngột trong lập luận hiện tại.

Vật chất tối không làm thay đổi nhiều logic cơ bản. Trong đa số lý thuyết, nó tương tác với vật chất thông thường yếu đến mức, đối với câu hỏi này, về cơ bản nó gần như đứng ngoài cuộc.

\begin{classroomqa}
\lecturetimestamp{01:24:49}
\audiencequestion{Hấp dẫn ảnh hưởng đến toàn bộ chuyện này như thế nào?}
\lecturerresponse{Chủ yếu theo hai cách. Thứ nhất, nếu không có sự giãn nở vũ trụ thì sẽ không có giai đoạn ngoài cân bằng. Thứ hai, sự bay hơi của lỗ đen cho một lý do mạnh để nghi ngờ rằng số baryon không thể là một định luật bảo toàn chính xác mãi mãi. Ngoài ra, các tương tác hấp dẫn thông thường giữa các hạt quá yếu để có ý nghĩa lớn ở đây.}
\end{classroomqa}

Tiếp theo đó có một câu hỏi về điện tích. Liệu cùng lập luận về lỗ đen ấy có đe dọa cả sự bảo toàn điện tích hay không? Câu trả lời là không. Một lỗ đen tích điện mang các đường sức điện vươn ra ngoài, và các đường ấy giữ dấu vết của điện tích. Một lỗ đen như vậy không bay hơi hết hoàn toàn; thay vào đó nó tiến tới một tàn dư cực trị. Vì thế lập luận về lỗ đen đe dọa số baryon, chứ không đe dọa điện tích.

Một câu hỏi khác hỏi liệu hỗn hợp quark cực đậm đặc ở thuở đầu có nên đơn giản được coi là một lỗ đen hay không. Câu trả lời là không. Vấn đề không phải là một khối đặc tĩnh. Vũ trụ khi đó đang giãn nở. Câu chuyện về lỗ đen là một thí nghiệm tư duy về số phận của số baryon, chứ không phải một mô tả theo nghĩa đen về vũ trụ sơ khai như một vùng bị bẫy.

\begin{classroomqa}
\lecturetimestamp{01:31:14}
\audiencequestion{Liệu ta có biết cách sửa đổi mô hình chuẩn để bao gồm những phản ứng như vậy không?}
\lecturerresponse{Khó khăn lại nằm ở phía ngược lại. Rất dễ thêm quá nhiều sự vi phạm số baryon và làm cho proton phân rã quá nhanh. Vấn đề thực sự là giữ cho proton bền đúng như mức hiện có, đồng thời vẫn cho phép các quá trình trong vũ trụ sơ khai mà ta mong muốn.}
\end{classroomqa}

\section{Entropy và các câu hỏi về neutron ở phần sau}

Quay lại việc kiểm đếm. Trong một vũ trụ giãn nở đoạn nhiệt, entropy được bảo toàn xấp xỉ. Đối với một hỗn hợp nóng của các hạt sơ cấp, entropy về cơ bản chính là số hạt, sai khác bởi các hệ số cỡ 1. Đó là lý do cuộc thảo luận không cần đến một sự tạo ra photon khổng lồ ở giai đoạn muộn từ hư không. Các quá trình hủy cặp chủ yếu chỉ tái phân bố thành phần của hỗn hợp thành bức xạ.

Phiên bản cũ của vấn đề đã kể câu chuyện theo chiều ngược lại. Thay vì hỏi vì sao có quá ít proton sống sót so với photon, người ta lại hỏi theo cách kia: vì sao có quá nhiều photon so với proton? Nhưng câu hỏi sắc hơn không phải là vì sao vũ trụ lại tạo ra nhiều entropy đến thế ở giai đoạn muộn. Mà là vì sao ngay từ đầu bất đối xứng giữa hạt và phản hạt lại nhỏ đến thế: chỉ khoảng 1 proton dư ra trên mỗi tỷ cặp proton--phản proton.

Một neutron tự do quả thật có phân rã, nhưng quá trình đó \emph{không} vi phạm số baryon:
\[
n \to p + e^- + \bar\nu.
\]
Neutron mang cùng số baryon như proton.

Đối với nơtron tự do, năng lượng phân rã khả dụng là rất nhỏ, vì khối lượng của nơtron chỉ nhỉnh hơn một chút so với khối lượng của trạng thái cuối cùng. Đó là một phần lý do khiến sự phân rã này diễn ra rất chậm theo tiêu chuẩn của các hạt cơ bản. Khi nơtron bị liên kết trong một hạt nhân, năng lượng liên kết có thể lấy đi hoàn toàn lượng năng lượng khả dụng nhỏ đó, và khi ấy nơtron liên kết có thể ổn định.

\begin{classroomqa}
\lecturetimestamp{01:43:03}
\audiencequestion{Các nơtron chỉ đơn giản là đang biến mất khỏi vũ trụ, hay chúng cũng được tạo ra?}
\lecturerresponse{Câu trả lời đầu tiên là các nơtron tự do bị mất đi do phân rã. Nhưng đó không phải toàn bộ câu chuyện. Có những nguồn thực sự tạo ra nơtron, bao gồm các phản ứng hạt nhân trong các sao và sự hình thành các sao nơtron. Vì vậy, tổng số nơtron không đơn giản chỉ là sự suy giảm một chiều.}
\end{classroomqa}